\section{Background and Related Work}

\subsection{LLMs as Behavioural Experimental Subjects and Computational Models}

Research on LLMs in the behavioural sciences has followed two broad approaches: using models as research tools to simulate participants, and treating models themselves as experimental subjects whose choices and context-dependent responses can be examined using psychological tasks \cite{Dillion2023ReplaceParticipants,Hagendorff2023MachinePsychology,Rahwan2019MachineBehaviour}. The former approach requires the model to support valid inferences about human populations, whereas the latter focuses on the measurable behaviour of a specified computational system. The validation criteria appropriate to these two approaches should not be conflated.

The present study adopts the second approach and defines an LLM as a computational model capable of producing comparable behaviour in cognitive tasks. This is an operational, rather than mechanistic, definition: it does not imply that the model shares human perceptual, learning, or decision-making processes. Research on machine psychology and machine behaviour similarly uses experimental paradigms to characterise artificial systems, but human-like content effects or aggregate behavioural fit may nevertheless arise from different underlying mechanisms \cite{Rahwan2019MachineBehaviour,Lampinen2024ContentEffects,Lin2025SixFallacies}.

This definition also limits the interpretation of LLMs as synthetic participants. Agent architectures incorporating memory, reflection, and planning can produce coherent patterns of behaviour, but those patterns depend on the design of the wider system \cite{Park2023GenerativeAgents}. Synthetic responses may also compress population heterogeneity or lose their original correspondence with human data when the model or prompt is changed \cite{Dillion2023ReplaceParticipants,Bisbee2024SyntheticSurveyData}. Accordingly, the present study examines the behavioural system jointly produced by the model, task, and measurement procedure; it does not treat LLM runs as substitutes for a human population.

\subsection{Reliability in the Behavioural Measurement of LLMs}

Conventional LLM evaluations often emphasise task accuracy, whereas behavioural experiments must also consider choice distributions, behavioural trajectories, and the extent to which an observed experimental effect depends on a particular measurement formulation. Evaluations using multiple templates indicate that reliance on a single prompt may overstate the certainty of a conclusion, while scoring procedures can themselves introduce apparent variation \cite{Sclar2024PromptFormatting,Mizrahi2024MultiPrompt,Hua2025FlawArtifact}. In this study, \emph{behavioural-measurement reliability} is operationalised as the degree to which the direction and magnitude of a behavioural conclusion remain stable across predefined prompt formulations when the model snapshot, task implementation, generation temperature, response format, and parser are held fixed. This definition does not encompass long-term test--retest reliability across service versions, external validity across other tasks, or the reliability of claims about underlying mechanisms.

The measurement procedure for an LLM includes the prompt, task interface, visible history, output constraints, parser, generation temperature, and deployment state. These components may alter the option produced by the model or affect how its response is parsed and scored \cite{Webson2022PromptMeaning,Sclar2024PromptFormatting,Loya2023PromptSensitivityDecision,Chen2024ChatGPTBehavior,Hua2025FlawArtifact}. A fixed action space and parser can reduce scoring ambiguity, but cannot eliminate the influence of prompt length, informational salience, or information position. Reliable interpretation therefore requires a clear distinction between the elements held constant, those deliberately manipulated, and the scope of the resulting inference.

The unit of analysis is equally important. Repeated runs share model parameters and deployment conditions; variation among them arises from their contexts, task environments, and generation stochasticity. Treating such runs as independent synthetic persons would overstate their population representativeness \cite{Dillion2023ReplaceParticipants,Bisbee2024SyntheticSurveyData}. The present study aligns complete LLM runs with participant-level summaries. It neither treats trials within a run as independent participants nor attributes genuine individual identities to repeated runs.

The coexistence of human--model similarities and differences documented by Lampinen et al.~\cite{Lampinen2024ContentEffects}, together with Lin's distinction between output similarity and mechanistic equivalence \cite{Lin2025SixFallacies}, demonstrates why a single judgement of being ``human-like'' is insufficient for evaluating LLM behavioural evidence. The present study consequently distinguishes three non-interchangeable analytical dimensions: prompt stability, task-specific behavioural outcomes, and human-reference proximity. Stability does not imply that a behaviour is preferable or more human-like, and proximity between means does not establish similarity in distributions or mechanisms. Conclusions should therefore be based on prespecified effect magnitudes, uncertainty, and patterns across metrics, rather than on a single aggregate score or an isolated comparison.

\subsection{Prompt Sensitivity and Behavioural Variation}

Research on prompt sensitivity has shown that label bias, demonstration order, formatting, and option order can all alter model outputs \cite{Zhao2021Calibrate,Lu2022PromptOrder,Sclar2024PromptFormatting,Pezeshkpour2024OptionOrder}. The observation that irrelevant or misleading instructions can still produce high task scores further indicates that successful performance does not necessarily demonstrate a stable understanding of the task \cite{Webson2022PromptMeaning}. The prompt must therefore be regarded as part of the behavioural measurement procedure.

The field has accordingly begun to move from searching for a single ``best prompt'' towards multi-prompt evaluation and instance-level quantification of sensitivity \cite{Mizrahi2024MultiPrompt,Zhuo2024ProSA}. Much of this evidence, however, comes from static benchmarks or single-round task responses. Loya et al.~\cite{Loya2023PromptSensitivityDecision} found that choices in the Horizon Task varied with prompts and generation temperature across three OpenAI models, but examined only one sequential decision paradigm and did not include a multilingual comparison.

The interpretation of a prompt effect depends on the type of manipulation involved. Formatting and label-order manipulations primarily test sensitivity to incidental design features, whereas role-play prompting has also been shown to alter model performance across several reasoning benchmarks \cite{Kong2024RolePlayPrompting}. In the present study, Instruction specificity, Role framing, and Task-specific construct emphasis are predefined as substantively meaningful manipulations that may affect rule representation, role context, or conceptual salience. Their behavioural consequences in sequential decision tasks remain empirical questions. Because these formulations also differ in length, keywords, and information position, the complete formulation is the object of inference. These correlated presentational differences constitute a design limitation and cannot be decomposed retrospectively from the observed condition effect. Consequently, a condition effect should not automatically be characterised as unreliability caused by irrelevant wording \cite{Sclar2024PromptFormatting,Hua2025FlawArtifact}.

\subsection{Model Variation and Reproducibility in LLM Behavioural Measurement}

Behavioural measurements are not necessarily invariant across models. Multi-model evaluations indicate that both baseline performance and prompt sensitivity may differ among models \cite{Sclar2024PromptFormatting,Mizrahi2024MultiPrompt,Zhuo2024ProSA}. A comparison of model means in a single condition cannot distinguish an overall behavioural difference from differential formulation sensitivity. It is therefore necessary first to estimate the condition-minus-baseline effect for each model and then to compare those effects directly.

Cross-model comparisons are also affected by changes in model services. Proprietary services and model versions may change over time, often through opaque update processes, and the same model name does not guarantee identical behaviour at different collection dates \cite{Chen2024ChatGPTBehavior}. Palmer, Smith and Spirling~\cite{Palmer2024ProprietaryModels} accordingly emphasise the transparency and reproducibility constraints associated with proprietary language models. The present study records requested and resolved model identifiers, collection dates, generation settings, prompt versions, and failure logs. Using fixed model snapshots can reduce, but cannot eliminate, temporal drift.

Taken together, this evidence indicates that cross-model behavioural comparisons should examine both baseline behaviour and formulation sensitivity, while restricting inference to the model versions actually tested \cite{Chen2024ChatGPTBehavior,Palmer2024ProprietaryModels}. Differences should not be interpreted as causal effects of architecture or scale. The models, reference directions, and generation settings used in the present study are specified in the Method section.

\subsection{Language Effects in LLM Behavioural Measurement}

The multilingual capabilities of an LLM are not proportional replications of its capabilities in English. Cross-language benchmarks show that accuracy and performance on culturally situated knowledge vary across languages and task levels. Moreover, improvements in a larger model's average multilingual performance do not necessarily produce greater cross-lingual consistency \cite{Zhang2023M3Exam,Qi2023CrossLingualConsistency}. Support for a language therefore does not imply that the same behavioural experiment has identical measurement properties in that language, and findings obtained in English cannot automatically be generalised to Chinese or Spanish.

Existing multilingual research has primarily used translated benchmarks, cross-lingual prompting, or multilingual chain-of-thought methods. Instruction tuning in English can transfer to other languages, but the degree of transfer depends on factors including pre-training coverage, task format, the language of demonstrations, and the reasoning strategy employed \cite{Muennighoff2023CrosslingualGeneralization,Shi2023MultilingualCoT}. Performance may also vary according to how instructions, contexts, examples, and output components are translated \cite{Mondshine2025BeyondEnglish}. These approaches have less often examined both the baseline behaviour and within-language prompt effects of a fixed model completing full behavioural task runs.

Such cross-language differences also make prompt translation an experimental operation in its own right. Semantic alignment does not guarantee identical measurement properties, because translation quality, linguistic distance from English, and pre-training coverage may all affect model behaviour \cite{Mondshine2025BeyondEnglish}. The present study therefore requires task-relevant informational equivalence rather than complete semantic or pragmatic equivalence across language versions. Observed differences are interpreted cautiously as language-associated differences specific to the tested model, translation, and task implementation. The translation review procedure and statistical contrasts are specified in the Method section.

\subsection{Sequential Decision Tasks and Their Cognitive Constructs}

In each of the three tasks selected for this study, a current choice and the feedback subsequently presented contribute jointly to a behavioural trajectory across trials or steps. This structure makes it possible to examine how the model uses previous information and adjusts later choices \cite{Wilson2014,Bechara1994,Lejuez2002,Loya2023PromptSensitivityDecision}. Because of this trajectory dependence, the complete run is used as the analysis unit; trials within the same run are not treated as independent participants.

The Horizon Task uses information asymmetry and the number of remaining decision opportunities to examine the explore--exploit dilemma \cite{Wilson2014}. Directed and random exploration can produce superficially similar choices while corresponding to different computational descriptions \cite{Gershman2018Exploration}. The information-seeking choice rate, horizon-related exploration change, and model-derived decision-noise difference should therefore be interpreted separately.

The IGT was originally developed to examine how repeated reward and loss feedback shapes subsequent choice \cite{Bechara1994}. However, the same eventual deck preference may arise from different patterns of outcome valuation, learning, and response consistency \cite{Busemeyer2002PVL,Steingroever2013PVLDelta}. Reviews of its construct validity likewise caution against treating an aggregate IGT score as a unitary measure of decision-making ability \cite{Buelow2009IGTValidity}. For an LLM, the explicit presentation of cumulative or average outcomes may impose a lower memory demand than the classic human task. Interpretation must therefore remain limited to the behavioural metrics aligned across implementations.

The BART examines sequential risk-taking and stopping behaviour by increasing potential gains while simultaneously increasing the risk of an explosion \cite{Lejuez2002}. The number of pumps is not a pure measure of risk preference; it may also reflect learning about explosion probabilities and response consistency \cite{Wallsten2005BARTModel,VanRavenzwaaij2011BART}. The present study therefore reports adjusted pumps, explosion rate, and post-explosion adjustment as complementary indicators. Although the three tasks address exploration, feedback-based learning, and risk-taking, respectively, they differ in their constructs, measurement scales, and feedback structures. Each task and metric is consequently analysed separately, and no aggregate score of general ``decision-making ability'' is constructed.

\subsection{Human-Reference Comparisons: Purpose and Limits}

Human-reference comparisons can describe where a model lies relative to an empirical human reference on a particular metric. The present analysis distinguishes the position of the model mean, the proportion of runs falling within the empirical human reference interval, and the response to an experimental manipulation. These quantities answer different questions and are not substitutes for one another. Humans and models may exhibit a similar aggregate effect while differing in task-specific behaviour \cite{Lampinen2024ContentEffects}, and correspondence at the output level does not demonstrate mechanistic equivalence \cite{Lin2025SixFallacies}. Human-reference proximity is therefore treated as a descriptive behavioural reference, not as a test of cognitive equivalence or shared mechanisms.

These comparisons are also constrained by the units of analysis and the degree of task alignment. An LLM run can be operationally aligned with a participant-level summary, but it is not an independent human individual; repeated runs share model parameters and deployment conditions \cite{Dillion2023ReplaceParticipants,Bisbee2024SyntheticSurveyData}. Differences in task length, information presentation, and familiarisation are treated as design-alignment limitations, and comparisons are restricted to metrics that can be aligned across the LLM and human implementations \cite{Wilson2014,Buelow2009IGTValidity,Wallsten2005BARTModel}. A mean deviation and empirical reference-interval coverage describe, respectively, relative mean position and the proportion of LLM runs within a prespecified human interval; they need not support the same descriptive conclusion.

The human reference is not a normative performance target. The direction of the IGT advantageous-choice rate is comparatively interpretable, but neither the number of BART pumps nor information seeking has a single context-independent optimum. Greater proximity to the human mean therefore does not automatically indicate improved decision quality \cite{Wilson2014,Buelow2009IGTValidity,Wallsten2005BARTModel}. The signed human-SD-scaled mean deviation, its absolute value, and empirical interval coverage are consequently reported as complementary supplementary or descriptive measures. They are not treated as Cohen's \(d\), evidence of overall equivalence, or evidence of a shared mechanism.

\subsection{Research Gap and Positioning of the Present Study}

Existing research has established several relevant lines of inquiry, including synthetic-participant simulation, machine psychology, prompt-robustness evaluation, and multilingual benchmarking. These approaches have generally pursued different objectives, such as population simulation, performance on individual benchmark items, the reproduction of particular psychological effects, or aggregate accuracy \cite{Dillion2023ReplaceParticipants,Hagendorff2023MachinePsychology,Mizrahi2024MultiPrompt,Zhang2023M3Exam}. Collectively, they demonstrate that LLM behaviour can be investigated experimentally, but they do not integrate measurement stability, task-specific behaviour, and human-reference proximity within a common evaluation design.

Some existing studies include multiple models or multiple prompt variants, but their evidence is drawn mainly from static benchmarks or from a single sequential decision task. Studies of multilingual consistency, meanwhile, typically use separate factual-knowledge or benchmark settings \cite{Loya2023PromptSensitivityDecision,Sclar2024PromptFormatting,Qi2023CrossLingualConsistency}. Within the literature reviewed here, no study was identified that combines, within a single design, full sequential task runs, controlled prompt variants, multiple models, multiple decision constructs, and a cross-language comparison holding the model fixed. The available evidence is therefore insufficient to establish whether behavioural conclusions remain robust across these prompt formulations, model versions, and language implementations.

The present study combines controlled multi-prompt evaluation with full sequential task runs and extends this approach across tasks, models, and languages \cite{Loya2023PromptSensitivityDecision,Sclar2024PromptFormatting,Mizrahi2024MultiPrompt}. It also distinguishes formulation stability, task-specific behavioural outcomes, and human-reference proximity. Its contribution is not to determine whether LLMs are globally ``like humans'', but to assess systematically whether conclusions about their behaviour in cognitive tasks remain robust across predefined variations in prompts, models, and languages.
